\section{Introduction}
\label{sec:introduction}

Digital platforms have increasingly replaced traditional revenue streams across various industries, including music, where streaming platforms now account for $84\%$ of total revenues in the US \citep{riaa}. Users seeking to consume music can either search for particular songs or albums or browse thousands of playlists---i.e., curated lists of songs centered around a specific genre, mood, activity, event, or artist. Playlists have contributed to a significant expansion in music consumption on streaming services \citep{Datta_etal_MktScie}, accounting for about $30$\% share of listening time on Spotify. Another $30$\% of music consumption is from tracks saved in users' music libraries \citep[][]{spotifyf1filing}, which often get discovered by consumers via playlists. %In the U.S., $27\%$ of users reported listening to curated playlists in 2017 (Q1), relative to $55\%$ in 2019 (Q2).\footnote{See \href{https://www.statista.com/statistics/868658/curated-playlists-streaming-music-services/}{{https://www.statista.com/statistics/868658/curated-playlists-streaming-music-services/}}, accessed 30 October, 2023.} 
Many of the most popular playlists are designed and professionally managed by Spotify. % to attract followers and meet user needs. 
With 15,000 playlists, Spotify attracts $63$\% of platform-wide playlist followers \citep[e.g.,][]{Aguiar_Waldfogel_power}.   %Users can also share their playlists and, in doing so, generate followers for their libraries.
%
%%% Maybe we can remove this altogether?%%% 
%
%
Other playlists on the platform are either curated by users or professional editors---including major and independent music labels. %, independent labels, artists, radio stations, brands, and influencers. %---create playlists by adding or removing tracks at fixed intervals based on the perceived fit between tracks and the theme of a playlist (e.g., genre, mood).\footnote{See, e.g., \href{https://artists.spotify.com/blog/behind-the-playlists-r-and-b}{``Behind the Playlists: R\&B"}, accessed July 19, 2022.}
%In the case of Spotify playlists, labels and artists can pitch songs to playlist curators for inclusion via a dedicated portal (``Spotify for Artists").\footnote{Spotify pays royalties relative to a label's share of total streams on the platform \citep{loudandclear}.} 
Finally, Spotify features some playlists across various categories on the Spotify app’s Search Page. This changing directory of lists features mostly playlists designed internally by Spotify. 

Despite the growing importance of playlists to music consumption, not much is known about why users prefer some playlists over others. For example, users may follow playlists because they are prominently featured on the platform, giving them a shortcut to finding music they want to listen to from Spotify's vast assortment. Alternatively, users may follow playlists because they contain tracks by famous artists such as Taylor Swift (Universal), Adele (Sony), or Ed Sheeran (Warner). In our empirical analysis, we compare how adding tracks from popular major label artists to a playlist causally affects daily playlist followers in the short term (i.e., using variation within a month) versus the impact of the playlist being featured on the Search Page of the Spotify app. Before our work, little was known about how strongly users respond to these drivers of playlist demand.\footnote{Prior related work has explored consumption patterns on digital media, including live streaming \citep[e.g.,][]{Live_streaming_MKTSCIE_21,jmr_livestreaming} and on-demand streaming services \citep[e.g.,][]{wlomert2023frontiers,Datta_etal_MktScie}. In addition, researchers have measured the market power of digital platforms \citep[e.g.,][]{ review_platform_power_23,economics_two_sided_markets,platform_power_Song_2021}. In the context of music streaming, for example, \cite{playlistbias} show that Spotify likely exercises its power to make curation decisions for the ``New Music Friday" playlist favor independent record labels and, to a lesser extent, female artists.}

We use two complementary methods to estimate these effects: a two-way fixed effects model \citep[e.g.,][]{shapiro_hitsch_tuchman_2021}, and a synthetic difference in difference model  
\citep[e.g.,][]{SDiD_AER_Arkhangelsky_21}. Our analysis builds on a unique data set collected from \textit{Chartmetric.com} that tracked over 30,000 playlists across 66 music categories between October 2019 and March 2020. The data includes the number of daily aggregate playlist followers as well as various characteristics such as whether playlists add tracks from famous artists of the three major record labels (Universal, Sony, and Warner), the share of content on the playlist owned by these labels, the duration of the playlist, and the acoustic attributes of tracks \citep{boughanmiacoustics}. Our data set also records which playlists were featured prominently on the platform's Search Page on a given day, obtained through live scraping from \textit{everynoise.com/worldbrowser.cgi} \citep{scraping}.

Our identification strategy relies on the synthetic difference in difference (SDID) method to estimate treatment effects (e.g., the effect of featuring playlists on the Search Page) relative to control playlists that closely mimic treated playlists but that have not received the treatment \citep[e.g.,][]{berman2022value}. For this purpose, the SDID method reweighs the large pool of control playlists to match pre-exposure trends between treated and control observations, reducing the reliance on parallel trend assumptions as in conventional difference in difference approaches \citep{SDiD_AER_Arkhangelsky_21}. We also infer treatment effects in a flexible two-way fixed effects (2W-FE) model accounting for date and playlist-month fixed effects \citep[e.g.,][]{shapiro_hitsch_tuchman_2021}. We find that the treatment effects are consistent between the SDID and 2W-FE models, demonstrating robustness with regard to different model specifications and assumptions. 

Our SDID estimates reveal that users strongly value playlists featured in the Spotify app, implying that daily followers go up by 0.95\% if a playlist is included on the Search Page. For large playlists like ``Reggaeton Classic" (which has 1.9 million followers), being prominently featured on the Spotify app can yield 18,050 additional daily followers. This finding indicates that Spotify is effective in guiding consumer searches of its vast playlist assortment through what it features on its Search Page. In comparison, we find that users are about half as sensitive to the inclusion in a playlist of tracks by popular major label artists---adding at least one track by a major label superstar artist increases daily followers by 0.45\%. 

Our analysis of playlist demand has two main implications. First, our estimates show that playlist curators attract followers twice as easily when their list is featured on the Search Page as when they add a track by a major label superstar. %This insight is of interest to policymakers. %, who are trying to understand the new balance of power between digital platforms like Spotify and the suppliers of digital goods \citep{digitalsinglemarket}.
Historically, the major music labels controlled content distribution in the music industry and exerted influence over how music was advertised and sold offline \citep[][]{prey_2020,du_gay_and_negus_1994}. Today, however, our findings corroborate researchers’ observations that music streaming platforms have higher influence over users’ listening behavior compared to labels%: Industry executives and artists fear that Spotify could exert unfair market power through its access to user data and the ability to
---such as by promoting songs via playlists to the advantage of some producers over others \citep[][]{Aguiar_Waldfogel_power}. 

Second, our results imply that superstar artists remain a relevant asset for major music labels. This insight is of interest to record labels and music streaming platforms, who are disputing royalty rates in the US.\footnote{See case docket \textit{19-CRB-0005-WR} at the Copyright Royalty Board, available at \href{https://app.crb.gov/case/detail/19-CRB-0005-WR\%20\%282021-2025\%29}{https://app.crb.gov/}, accessed March 5, 2024.} Our findings could aid judges in assessing the consumer value of major label content on music streaming platforms: the inclusion of major label superstar tracks in playlists draws significant attention, as reflected by its positive (short-term) effect on followers. %These insights can be used to negotiate royalty rates between music labels and streaming platforms.   


%Some industry representatives started worrying about the value of their superstars as major labels' market share on streaming services have been declining. However, our analysis underscores the continued relevance of superstars in the digital age: the inclusion of major label superstar tracks in playlists draws significant attention, as reflected by its positive effect on daily followers. 

The rest of the paper is organized as follows. Section \ref{sec:data} presents the data. Section \ref{identification_strategy} discusses causal identification challenges and presents two model specifications that address these challenges: the two-way fixed effects model and the synthetic difference in difference model. Section \ref{results} presents the treatment effect estimates and discusses robustness checks. % 
We offer concluding thoughts and discuss the implications of our results in Section \ref{conclusion}.